\section{Σύνθεση της Υποδομής Ανάπτυξης}
\label{sec:appendix_compose_stack}

Το \texttt{\ENG{central-docker-infrastructure}} δηλώνει ένα σαφώς ορισμένο \ENG{docker} περιβάλλον,
το οποίο λειτουργεί ως κοινό πλαίσιο ανάπτυξης (\ENG{local-operational stack}),
μέσα στο οποίο ο \ENG{engine}, ο \ENG{player}, το \ENG{LMS}, τα άλλα διάφορα \ENG{services} 
συνεργάζονται πάνω στο κοινό δίκτυο \texttt{\ENG{scorm-network}}.

Η δήλωση των υπηρεσιών περιλαμβάνει τρία διαφορετικά επίπεδα ευθύνης. Το πρώτο είναι το \ENG{application layer},
δηλαδή τα \texttt{\ENG{engine}}, \texttt{\ENG{player}} και \texttt{\ENG{lms}}. Το δεύτερο είναι το \ENG{state and object layer},
όπου ανήκουν τα \texttt{\ENG{postgres}}, \texttt{\ENG{redis}}, \texttt{\ENG{minio}} και το \texttt{\ENG{minio-init}}.
Το τρίτο είναι το \ENG{observability layer}, όπου εντάσσονται τα \texttt{\ENG{logstash}}, \texttt{\ENG{elasticsearch}} 
και \texttt{\ENG{kibana}}. Η τριμερής αυτή διάκριση υπάρχει για να διευκολύνει την κατανόηση της σύνθεσης,
αλλά και για να αναδείξει τη λογική οργάνωσης των υπηρεσιών.


\begin{figure}[H]
  \centering
\begin{chapterfigurebox}
  \begin{tikzpicture}[
      font=\footnotesize,
      node distance=1.7cm and 2.1cm,
      box/.style={draw, rounded corners, fill=gray!5, align=center, text width=3.0cm, minimum height=1.05cm},
      emph/.style={draw, rounded corners, fill=gray!12, align=center, text width=3.3cm, minimum height=1.05cm, font=\footnotesize\bfseries},
      store/.style={draw, cylinder, shape border rotate=90, aspect=0.28, minimum height=1.15cm, minimum width=3.0cm, align=center, fill=gray!5},
      boundary/.style={draw, dashed, rounded corners, inner sep=11pt},
      arrow/.style={-Latex, thick}
    ]
    \node[emph] (lms) {\ENG{LMS}\\\texttt{\ENG{8000}}};
    \node[emph, right=of lms] (player) {\ENG{Player}\\\texttt{\ENG{3000}}};
    \node[emph, right=of player] (engine) {\ENG{Engine}\\\texttt{\ENG{8080}}};

    \node[store, below=1.8cm of lms] (postgres) {\ENG{Postgres}\\\texttt{\ENG{5432}}};
    \node[store, below=1.8cm of player] (redis) {\ENG{Redis}\\\texttt{\ENG{6379}}};
    \node[store, below=1.8cm of engine] (minio) {\ENG{MinIO}\\\texttt{\ENG{9000/9001}}};

    \node[box, below=1.8cm of minio] (minioinit) {\ENG{minio-init}};
    \node[box, right=2.4cm of engine] (logstash) {\ENG{Logstash}\\\texttt{\ENG{12201/udp}}};
    \node[box, below=0.9cm of logstash] (elastic) {\ENG{Elasticsearch}\\\texttt{\ENG{9200}}};
    \node[box, below=0.9cm of elastic] (kibana) {\ENG{Kibana}\\\texttt{\ENG{5601}}};

    \draw[arrow] (lms) -- (engine);
    \draw[arrow] (player) -- (engine);
    \draw[arrow] (engine) -- (postgres);
    \draw[arrow] (engine) -- (redis);
    \draw[arrow] (engine) -- (minio);
    \draw[arrow] (player) -- (minio);
    \draw[arrow] (minioinit) -- (minio);
    \draw[arrow] (engine) -- (logstash);
    \draw[arrow] (player) -- (logstash);
    \draw[arrow] (lms) -- (logstash);
    \draw[arrow] (logstash) -- (elastic);
    \draw[arrow] (elastic) -- (kibana);

    \node[boundary, fit=(lms)(player)(engine)(postgres)(redis)(minio)(minioinit)(logstash)(elastic)(kibana)] (compose) {};
    \node[font=\footnotesize, align=center, text width=6.4cm, above=0.16cm of compose] {\ENG{docker-compose stack on scorm-network}};
  \end{tikzpicture}
\end{chapterfigurebox}
  \caption{Η δηλωμένη σύνθεση υπηρεσιών της υποδομής ανάπτυξης, όπως προκύπτει από το \texttt{\ENG{docker-compose.yml}}.}
  \label{fig:appendix_compose_topology}
\end{figure}


\begin{table}[H]
\centering
\caption{Κύριες υπηρεσίες της σύνθεσης, \ENG{host ports} και λειτουργικός ρόλος}
\label{tab:appendix_service_port_store_matrix}
\begingroup
\footnotesize
\setlength{\tabcolsep}{3pt}
\renewcommand{\arraystretch}{1.20}
\begin{adjustbox}{center,max width=\textwidth}
\begin{tabular}{|p{2.4cm}|p{2.0cm}|p{4.6cm}|p{4.5cm}|}
\hline
Υπηρεσία & \ENG{Host port} & Κύριος ρόλος & Μορφή αποθήκευσης ή ελέγχου \\
\hline
\texttt{\ENG{lms}} & \texttt{\ENG{8000}} & \ENG{Laravel} \ENG{UI} και διοικητικός/μαθησιακός πελάτης & \ENG{SQLite} τοπικά, \ENG{bootstrap} μέσω \texttt{\ENG{entrypoint.sh}} \\
\hline
\texttt{\ENG{player}} & \texttt{\ENG{3000}} & \ENG{Launch shell}, \ENG{content serving}, εσωτερικά \ENG{commit/terminate routes} & \ENG{player-cache volume}, \ENG{health endpoint} \\
\hline
\texttt{\ENG{engine}} & \texttt{\ENG{8080}} & \ENG{SCORM Engine REST API}, \ENG{launch/runtime} πυρήνας & \ENG{health endpoint}, εξάρτηση από \ENG{Postgres/Redis/MinIO} \\
\hline
\texttt{\ENG{postgres}} & \texttt{\ENG{5432}} & Κύριο σχεσιακό αποθετήριο μόνιμης αποθήκευσης & \texttt{\ENG{postgres-data}} volume και \ENG{pg\_isready} \ENG{healthcheck} \\
\hline
\texttt{\ENG{redis}} & \texttt{\ENG{6379}} & Ενεργή \ENG{runtime} κατάσταση \ENG{launches} & \texttt{\ENG{redis-data}} volume και \ENG{redis-cli ping} \\
\hline
\texttt{\ENG{minio}} & \texttt{\ENG{9000/9001}} & \ENG{Object storage} για \ENG{course packages} & \texttt{\ENG{minio-data}} volume, \ENG{live health probe}, \ENG{bucket init} \\
\hline
\texttt{\ENG{logstash}} & \texttt{\ENG{12201/udp}} & \ENG{GELF log ingestion} & \ENG{pipeline config} από \texttt{\ENG{logstash.conf}} \\
\hline
\texttt{\ENG{elasticsearch}} & \texttt{\ENG{9200}} & \ENG{Log indexing} και αναζήτηση & \texttt{\ENG{es-data}} volume και εσωτερική \ENG{HTTP probe} \\
\hline
\texttt{\ENG{kibana}} & \texttt{\ENG{5601}} & \ENG{Log inspection UI} & \ENG{status endpoint} και σύνδεση με \ENG{Elasticsearch} \\
\hline
\texttt{\ENG{minio-init}} & --- & Μεταβατική δημιουργία \ENG{buckets} κατά την εκκίνηση & Επιτρεπτός τερματισμός με \ENG{exit code 0} στην \ENG{CI} λογική \\
\hline
\end{tabular}
\end{adjustbox}
\endgroup

\end{table}
